%\begin{center}
%\large \bf \runtitulo
%\end{center}
%\vspace{1cm}
\chapter*{\runtitulo}

\noindent 
Los transformers son el componente básico de los grandes modelos de lenguaje (LLM).
En esta tesis, estudiamos los lenguajes formales que reconocen las diferentes abstracciones teóricas de los transformers cuando se limita el número de pasos de chain-of-thought.
En primer lugar, partimos del trabajo de \citet{toyota} y estudiamos las propiedades de su abstracción.
En segundo lugar, ampliamos su modelo para incorporar el comportamiento estocástico de los transformers, al que se le ha prestado poca atención en la literatura.
Estudiamos los lenguajes formales reconocibles por nuestra definición de transformers probabilísticos cuando se limita el número de pasos de chain-of-thought.
Nuestro análisis se lleva a cabo desde la perspectiva de la teoría de circuitos booleanos.
Finalmente, relacionamos ambos modelos con las máquinas de Turing y ofrecemos argumentos sobre cuándo se puede y cuándo no se puede ignorar la aleatoriedad de los transformers para estudiar su poder expresivo.

Nuestros principales resultados indican que los transformers que realizan $T(n)$ pasos de chain-of-thought son equivalentes a circuitos booleanos de tamaño $T(n)$.
Esto es válido tanto para los transformers y circuitos determinísticos como para los probabilísticos, incluso en el contexto $\Ptime$-uniforme o no uniforme.

\bigskip

\noindent\textbf{Palabras claves:} transformers probabilísticos, chain-of-thought, circuitos booleanos, lenguajes formales.
